% =============================================================
%  ch09_pcb.tex  --  Optimierungspotenzial für ein eigenes PCB
% =============================================================

\chapter{Optimierungspotenzial für ein eigenes PCB}

Diese Punkte sind \emph{Erweiterung}, nicht Voraussetzung -- mit Modulen kommst
du zunächst aus. Für ein späteres eigenes Board:
\begin{itemize}
\item \textbf{Treiber integrieren} statt Stepstick-Module: TMC2209 direkt
  bestücken -- spart Bauhöhe, bessere Thermik über Kupferflächen und Thermovias
  unter dem Chip.
\item \textbf{Strom-/Spannungsmessung onboard:} INA-Sensor mit dimensioniertem
  Shunt je Kanal nahe am Treiber, kurze Kelvin-Anbindung des Shunts.
\item \textbf{Power-Integrity:} Bulk-Elko plus keramische Stützkondensatoren
  direkt an jedem Treiber; getrennte Power-/Logik-Masse mit Sternpunkt
  (Single-Point-Ground).
\item \textbf{Signalführung:} STEP/DIR bzw.\ UART als kurze, von der Leistung
  getrennte Leitungen; bei gemischten Pegeln (3{,}3\,V/5\,V) Levelshifter
  vorsehen.
\item \textbf{Schutz:} ESD-Dioden an USB- und Encoder-Leitungen; Verpolschutz;
  Freilauf-/EMV-Betrachtung der Motorleitungen.
\item \textbf{Konnektivität:} JST-Stecker für Motoren/Encoder, Schraubklemme
  für die Versorgung; optional CAN-Transceiver für saubere
  Mehrachs-Verkettung statt vieler UARTs.
\item \textbf{Onboard-Encoder-Interface} für späteren Closed-Loop-Betrieb (FOC).
\item \textbf{Thermik:} Thermopad/Kühlfläche unter den Treibern, definierter
  Luftstrom.
\end{itemize}
